% --- CORRECCIÓN IMPORTANTE EN LA PRIMERA LÍNEA ---
% Agregamos "xcolor=table" para que funcione \rowcolor sin errores
\documentclass[aspectratio=169, 10pt, xcolor=table]{beamer}

% --- Configuración del Tema (estilo profesional como el ejemplo) ---
\usetheme{Madrid}
\usecolortheme{dolphin}

% --- Paquetes necesarios ---
\usepackage[utf8]{inputenc}
\usepackage[spanish]{babel}
\usepackage{amsmath, amssymb}
\usepackage{booktabs}
\usepackage{graphicx}
\usepackage{listings}
\usepackage{tikz}
\usetikzlibrary{arrows.meta, positioning, shapes.geometric}

% --- Ruta de Imágenes (crea la carpeta "imagenes" y coloca ahí tus archivos) ---
\graphicspath{{./imagenes/}}

% --- Estilo para código Python ---
\definecolor{codegreen}{rgb}{0,0.6,0}
\definecolor{codegray}{rgb}{0.5,0.5,0.5}
\definecolor{codepurple}{rgb}{0.58,0,0.82}
\definecolor{backcolour}{rgb}{0.95,0.95,0.92}
\lstset{
    language=Python,
    backgroundcolor=\color{backcolour},
    commentstyle=\color{codegreen},
    keywordstyle=\color{blue},
    numberstyle=\tiny\color{codegray},
    stringstyle=\color{codepurple},
    basicstyle=\ttfamily\scriptsize,
    breaklines=true,
    frame=single,
    showstringspaces=false
}

% --- Pie de página personalizado ---
\makeatletter

\setbeamertemplate{footline}{
  \leavevmode%
  \hbox{%
  \begin{beamercolorbox}[wd=.333333\paperwidth,ht=2.25ex,dp=1ex,center]{author in head/foot}%
    \usebeamerfont{author in head/foot}\insertshortauthor
  \end{beamercolorbox}%
  \begin{beamercolorbox}[wd=.333333\paperwidth,ht=2.25ex,dp=1ex,center]{title in head/foot}%
    \usebeamerfont{title in head/foot}Sesión 3
  \end{beamercolorbox}%
  \begin{beamercolorbox}[wd=.333333\paperwidth,ht=2.25ex,dp=1ex,right]{date in head/foot}%
    \usebeamerfont{date in head/foot}NLP \hfill \insertframenumber/\inserttotalframenumber \hspace*{0.3cm}
  \end{beamercolorbox}}%
  \vskip0pt%
}

\makeatother

% --- Datos de la portada ---
\title[SESIÓN 3 NLP]{\textbf{Sesión 3}}
\subtitle{Representación del Texto: \\ BOW, TF-IDF y N-gramas}
\author[Prof. C. Sánchez]{Carlos César Sánchez Coronel}
\institute{\textbf{NLP}}
\date{}

% Logo opcional (descomentar si tienes la imagen)
% \titlegraphic{\includegraphics[height=1.5cm]{logo.png}}

% --- Eliminar la barra de navegación (opcional) ---
\setbeamertemplate{navigation symbols}{}

% --- Índice al inicio de cada sección ---
\AtBeginSection[]{
  \begin{frame}
    \frametitle{Contenido}
    \tableofcontents[currentsection]
  \end{frame}
}

% ============================================================================
% INICIO DEL DOCUMENTO
% ============================================================================
\begin{document}

% --- Portada ---
\begin{frame}
    \titlepage
\end{frame}

% --- Índice general ---
\begin{frame}{Contenido}
    \tableofcontents
\end{frame}

% ============================================================================
% SECCIÓN 1: INTRODUCCIÓN
% ============================================================================
\section{Introducción}

\begin{frame}{¿Por qué representar el texto numéricamente?}
    \begin{itemize}
        \item Los modelos de machine learning operan con números, no con texto directamente.
        \item Cada documento debe convertirse en un \textbf{vector numérico}.
        \item La representación elegida afecta:
        \begin{itemize}
            \item Capacidad semántica.
            \item Dimensionalidad.
            \item Rendimiento de los modelos.
        \end{itemize}
        \item En esta sesión: representaciones \textbf{dispersas} clásicas.
    \end{itemize}
    % Basado en introducción del PDF
\end{frame}

% ============================================================================
% SECCIÓN 2: BAG OF WORDS (BOW)
% ============================================================================
\section{Bag of Words (BOW)}

\begin{frame}{Definición e intuición}
    \begin{block}{Idea}
        Representar cada documento como una \textbf{bolsa} de sus palabras, ignorando el orden pero manteniendo la frecuencia.
    \end{block}
    \begin{itemize}
        \item El contenido se infiere por las palabras que contiene.
        \item \textbf{Vocabulario}: conjunto de todas las palabras únicas del corpus.
    \end{itemize}
    % Basado en PDF
\end{frame}

\begin{frame}{Construcción paso a paso}
    \begin{enumerate}
        \item Crear vocabulario con todas las palabras únicas del corpus.
        \item Asignar un índice a cada palabra.
        \item Construir la matriz documento-término: cada fila = documento, cada columna = palabra, valor = frecuencia.
    \end{enumerate}
    \vspace{0.3cm}
    \begin{exampleblock}{Ejemplo (3 documentos)}
        d1: "el gato persigue al ratón" \\
        d2: "el perro persigue al gato" \\
        d3: "el ratón huye del gato"
    \end{exampleblock}
    % Basado en PDF
\end{frame}

\begin{frame}{Ejemplo numérico de BOW}
    \textbf{Vocabulario}: al, el, gato, huye, del, perro, persigue, ratón \\
    \textbf{Matriz de frecuencias}:
    \[
    \mathbf{X} = \begin{pmatrix}
    d_1: & [1, 1, 1, 0, 0, 0, 1, 1] \\
    d_2: & [1, 1, 1, 0, 0, 1, 1, 0] \\
    d_3: & [0, 1, 1, 1, 1, 0, 0, 1]
    \end{pmatrix}
    \]
    % Basado en PDF
\end{frame}

\begin{frame}{Limitaciones de BOW}
    \begin{itemize}
        \item \textbf{Pérdida de orden}: "el gato persigue al ratón" vs "el ratón persigue al gato" tienen el mismo vector.
        \item \textbf{Alta dimensionalidad}: vocabulario enorme, vectores muy dispersos.
        \item \textbf{Falta de semántica}: sinónimos (gato/felino) son independientes.
        \item \textbf{Palabras comunes dominan}: "el", "de" aparecen en todos.
    \end{itemize}
    % Basado en PDF
\end{frame}

% ============================================================================
% SECCIÓN 3: N-GRAMAS
% ============================================================================
\section{N-gramas}

\begin{frame}{Definición de n-gramas}
    \begin{block}{N-grama}
        Secuencia contigua de $n$ ítems (palabras o caracteres) extraída de un texto.
    \end{block}
    \begin{itemize}
        \item \textbf{Unigramas} ($n=1$): palabras individuales (BOW).
        \item \textbf{Bigramas} ($n=2$): pares de palabras consecutivas.
        \item \textbf{Trigramas} ($n=3$): tripletes.
        \item También n-gramas de caracteres.
    \end{itemize}
    % Basado en PDF
\end{frame}

\begin{frame}{Ejemplo con bigramas}
    Frase: "el gato persigue al ratón"
    \begin{itemize}
        \item Bigramas: "el gato", "gato persigue", "persigue al", "al ratón".
        \item Capturan contexto local que los unigramas ignoran.
        \item Útil para detectar frases hechas o patrones.
    \end{itemize}
    % Basado en PDF
\end{frame}

\begin{frame}{Compensación: orden vs dimensionalidad}
    \begin{itemize}
        \item Los n-gramas aumentan drásticamente la dimensionalidad.
        \item Para un vocabulario de $m$ palabras, hay hasta $m^n$ n-gramas posibles.
        \item En la práctica, muchos no aparecen y se pueden filtrar.
        \item Compromiso: capturan algo de orden sin modelar estructura completa.
    \end{itemize}
    % Basado en PDF
\end{frame}

% ============================================================================
% SECCIÓN 4: TF-IDF
% ============================================================================
\section{TF-IDF}

\begin{frame}{Definición e intuición de TF-IDF}
    \begin{block}{TF-IDF (Term Frequency - Inverse Document Frequency)}
        Ponderación que da mayor peso a términos frecuentes en un documento pero raros en el corpus.
    \end{block}
    \begin{itemize}
        \item Penaliza palabras comunes (artículos, preposiciones).
        \item Resalta términos distintivos del documento.
    \end{itemize}
    % Basado en PDF
\end{frame}

\begin{frame}{Componentes de TF-IDF}
    \begin{columns}[t]
        \column{0.5\textwidth}
        \textbf{TF (Term Frequency)}
        \begin{itemize}
            \item Frecuencia bruta: $tf(t,d) = f_{t,d}$
            \item Logarítmica: $tf(t,d) = \log(1 + f_{t,d})$
            \item Normalizada: $tf(t,d) = \frac{f_{t,d}}{\sum_{t'} f_{t',d}}$
        \end{itemize}
        \column{0.5\textwidth}
        \textbf{IDF (Inverse Document Frequency)}
        \[
        idf(t) = \log\left(\frac{N}{df(t)}\right)
        \]
        \begin{itemize}
            \item $N$: número total de documentos.
            \item $df(t)$: número de documentos que contienen $t$.
            \item Suavizado: $\log\left(\frac{N+1}{df(t)+1}\right) + 1$
        \end{itemize}
    \end{columns}
    \vspace{0.3cm}
    \textbf{Peso TF-IDF}: $tfidf(t,d) = tf(t,d) \times idf(t)$
    % Basado en PDF
\end{frame}

\begin{frame}{Ejemplo numérico de TF-IDF}
    Retomando el corpus anterior:
    \begin{itemize}
        \item $N=3$, término "gato": $df=3$, $idf = \log(3/3)=0$.
        \item Término "perro": $df=1$, $idf = \log(3/1) \approx 1.099$.
        \item Término "persigue": aparece en d1 y d2, $df=2$, $idf \approx 0.405$.
    \end{itemize}
    Matriz TF-IDF (TF bruta):
    \[
    \begin{pmatrix}
    d_1: & [0,0,0,0,0,0,0.405,0.405] \\
    d_2: & [0,0,0,0,0,1.099,0.405,0] \\
    d_3: & [0,0,0,1.099,1.099,0,0,0.405]
    \end{pmatrix}
    \]
    % Basado en PDF
\end{frame}

\begin{frame}{Normalización de vectores TF-IDF}
    \begin{itemize}
        \item Es común normalizar cada vector documento a norma unitaria (por ejemplo, L2).
        \item La similitud del coseno se convierte en producto punto.
        \[
        \mathbf{v}_{\text{norm}} = \frac{\mathbf{v}}{\|\mathbf{v}\|_2}
        \]
        \item Así, la longitud del documento no influye en la similitud.
    \end{itemize}
    % Basado en PDF
\end{frame}

\begin{frame}{Interpretación geométrica}
    \begin{itemize}
        \item Cada documento es un punto en un espacio de alta dimensión.
        \item La similitud entre documentos se mide con el \textbf{coseno del ángulo}:
        \[
        \text{sim}(d_i, d_j) = \frac{\mathbf{v}_i \cdot \mathbf{v}_j}{\|\mathbf{v}_i\| \|\mathbf{v}_j\|}
        \]
        \item Si los vectores están normalizados, es el producto punto.
        \item Coseno = 1: misma dirección, muy similares.
        \item Coseno = 0: ortogonales, no relacionados.
    \end{itemize}
    % Basado en PDF y Pregunta 6
\end{frame}

% ============================================================================
% SECCIÓN 5: COMPARACIÓN CON EMBEDDINGS
% ============================================================================
\section{Comparación con Embeddings}

\begin{frame}{BOW/TF-IDF vs Word Embeddings}
    \begin{table}
        \centering
        \begin{tabular}{l|c|c|c}
            \toprule
            \textbf{Característica} & \textbf{BOW} & \textbf{TF-IDF} & \textbf{Word Embeddings} \\
            \midrule
            Captura orden & No & No & No (en promedio) \\
            Captura semántica & No & No & Sí \\
            Peso de palabras comunes & Alto & Bajo & Depende \\
            Dimensionalidad & Alta (dispersa) & Alta (dispersa) & Baja (densa) \\
            Interpretabilidad & Sí & Sí & Baja \\
            Requerimiento de datos & Bajo & Bajo & Alto \\
            Uso típico & Línea base, spam & Búsqueda, clasificación & NLP profundo \\
            \bottomrule
        \end{tabular}
    \end{table}
    % Basado en PDF
\end{frame}

\begin{frame}{¿Cuándo usar cada una?}
    \begin{itemize}
        \item \textbf{BOW}: modelos simples, línea base, interpretabilidad crucial.
        \item \textbf{TF-IDF}: cuando se necesita ponderar importancia, búsqueda de información, extracción de keywords.
        \item \textbf{Embeddings}: cuando hay suficientes datos y se requiere capturar semántica (sinónimos, analogías).
    \end{itemize}
    % Basado en PDF y Pregunta 8
\end{frame}

% ============================================================================
% SECCIÓN 6: LIMITACIONES DE REPRESENTACIONES DISPERSAS
% ============================================================================
\section{Limitaciones de Representaciones Dispersas}

\begin{frame}{Limitaciones}
    \begin{itemize}
        \item \textbf{Maldición de la dimensionalidad}: espacio enorme, datos muy dispersos.
        \item \textbf{Problemas de memoria y computación}: se requieren matrices sparse.
        \item \textbf{Falta de generalización semántica}: sinónimos son independientes.
        \item \textbf{Necesidad de reducción de dimensionalidad}: LSA (SVD) para obtener representaciones densas latentes.
    \end{itemize}
    % Basado en PDF
\end{frame}

% ============================================================================
% SECCIÓN 7: APLICACIONES PRÁCTICAS
% ============================================================================
\section{Aplicaciones Prácticas}

\begin{frame}{Aplicaciones de BOW y TF-IDF}
    \begin{itemize}
        \item \textbf{Clasificación de textos}: spam, noticias, sentimiento (con Naive Bayes, SVM).
        \item \textbf{Búsqueda de documentos similares}: similitud del coseno para recuperación.
        \item \textbf{Extracción de palabras clave}: términos con mayor TF-IDF en un documento.
    \end{itemize}
    % Basado en PDF
\end{frame}

% ============================================================================
% SECCIÓN 8: IMPLEMENTACIÓN EN PYTHON CON SCIKIT-LEARN
% ============================================================================
\section{Implementación en Python con scikit-learn}

\begin{frame}[fragile]{CountVectorizer (BOW)}
    \begin{lstlisting}
from sklearn.feature_extraction.text import CountVectorizer

corpus = [
    "el gato persigue al ratón",
    "el perro persigue al gato",
    "el ratón huye del gato"
]

vectorizer = CountVectorizer()
X = vectorizer.fit_transform(corpus)
print(vectorizer.get_feature_names_out())
print(X.toarray())
    \end{lstlisting}
\end{frame}

\begin{frame}[fragile]{TfidfVectorizer}
    \begin{lstlisting}
from sklearn.feature_extraction.text import TfidfVectorizer

vectorizer = TfidfVectorizer(norm='l2', use_idf=True, smooth_idf=True)
X = vectorizer.fit_transform(corpus)
print(vectorizer.get_feature_names_out())
print(X.toarray())  # valores TF-IDF normalizados L2
    \end{lstlisting}
\end{frame}

\begin{frame}[fragile]{Uso de n-gramas}
    \begin{lstlisting}
# Bigramas (unigramas y bigramas)
vectorizer = CountVectorizer(ngram_range=(1,2))
X = vectorizer.fit_transform(corpus)
print(vectorizer.get_feature_names_out())
    \end{lstlisting}
\end{frame}

\begin{frame}[fragile]{Ejemplo con dataset real (20 Newsgroups)}
    \begin{lstlisting}
from sklearn.datasets import fetch_20newsgroups
from sklearn.feature_extraction.text import TfidfVectorizer
from sklearn.naive_bayes import MultinomialNB
from sklearn.pipeline import make_pipeline

categories = ['alt.atheism', 'soc.religion.christian']
newsgroups = fetch_20newsgroups(categories=categories)

pipeline = make_pipeline(TfidfVectorizer(), MultinomialNB())
pipeline.fit(newsgroups.data, newsgroups.target)
    \end{lstlisting}
\end{frame}

% ============================================================================
% SECCIÓN 9: PREPARACIÓN PARA ENTREVISTAS
% ============================================================================
\section{Preparación para Entrevistas}

\begin{frame}{Preguntas típicas}
    \begin{itemize}
        \item Compara Bag of Words, TF-IDF y word embeddings. ¿Cuándo usarías cada uno?
        \item ¿Qué es IDF y por qué se usa?
        \item ¿Cómo calcularías la similitud entre dos documentos con TF-IDF?
        \item ¿Qué problemas tiene la representación BOW?
        \item ¿Qué son los n-gramas y para qué sirven?
    \end{itemize}
    % Basado en PDF y solucionario
\end{frame}

% ============================================================================
% SECCIÓN 10: CONEXIÓN CON EL RESTO DEL CURSO
% ============================================================================
\section{Conexión con el resto del curso}

\begin{frame}{Próximos pasos}
    \begin{itemize}
        \item \textbf{Sesión 4}: Clasificación de textos usando estas representaciones.
        \item \textbf{Sesión 5}: Word embeddings, que superan limitaciones semánticas.
        \item \textbf{Sesión 6}: Modelos secuenciales (RNN) que capturan orden.
        \item \textbf{Sesión 7}: Transformers, que usan subword tokenization y embeddings.
    \end{itemize}
\end{frame}

% ============================================================================
% SECCIÓN 11: RESUMEN
% ============================================================================
\section{Resumen}

\begin{frame}{Lo que hemos aprendido}
    \begin{exampleblock}{Conceptos clave}
        \begin{itemize}
            \item \textbf{Bag of Words (BOW)}: vector de frecuencias de términos. Simple, pero limitado (pérdida de orden, alta dimensionalidad, falta de semántica).
            \item \textbf{N-gramas}: extienden BOW capturando contexto local; aumentan dimensionalidad.
            \item \textbf{TF-IDF}: pondera términos según importancia en documento y rareza en corpus. Penaliza palabras comunes.
            \item Representaciones dispersas: sufren maldición de dimensionalidad, pero son interpretables y útiles para clasificación y búsqueda.
            \item Implementación con scikit-learn: \texttt{CountVectorizer}, \texttt{TfidfVectorizer}, n-gramas.
        \end{itemize}
    \end{exampleblock}
\end{frame}

% --- Diapositiva final ---
\begin{frame}
    \centering
    \Huge \textbf{¡Gracias!}

\end{frame}

\end{document}